\section{Related Work}

\subsection{Neural Product Retrieval}

Product search combines the vocabulary mismatch of Web retrieval with heterogeneous catalog fields and sparse behavioral labels. Lexical methods such as BM25 remain strong components, while neural dual encoders support semantic candidate generation and cross-encoders support more expensive reranking~\cite{robertson2009bm25,reimers2019sentencebert,karpukhin2020dpr,nogueira2019passage}. Hybrid and field-aware systems combine these signals. Choi et al. learn Transformer representations for multi-instance product fields and explicitly study structured matching together with lexical features~\cite{choi2020semantic}; related work extracts query attributes for semantic retrieval and ranking~\cite{loughnane2024attributes}. WANDS and the Amazon Shopping Queries Dataset provide public catalog fields and graded product-search judgments~\cite{chen2022wands,reddy2022shopping}. Dense embedding families and e-commerce-specific encoders broaden the available matching functions~\cite{xiao2024cpack,li2023gte}. These methods primarily improve matching effectiveness while generally treating the chosen product serialization as fixed.

\subsection{Content and Visibility Optimization}

A separate line of research studies how content can be changed to gain visibility. GEO evaluates presentational modifications intended to increase a source's presence in generative answers~\cite{aggarwal2024geo}. E-GEO specializes this objective to product descriptions and generative shopping rankings~\cite{bagga2025egeo}. SAGEO Arena restores retrieval and reranking to the evaluation boundary and shows why optimization must be assessed across the full search-augmented generation pipeline~\cite{kim2026sageo}. Earlier work likewise studies ranking-incentivized, quality-preserving content modification~\cite{goren2020ranking}.

Our objective differs. GEO asks how content can be modified to obtain more visibility. We ask whether semantically irrelevant representation choices should change visibility at all. The intervention is symmetric: a product can be helped or harmed, and both outcomes count as instability. We neither optimize a target merchant nor infer commercial benefit.

\subsection{Representation and Serialization Robustness}

Behavioral IR tests reveal model properties that aggregate effectiveness can hide. ABNIRML uses controlled probes for neural rankers~\cite{macavaney2022abnirml}; language-model studies show that word order can be simultaneously redundant for some predictions and consequential for others~\cite{hessel2021order,abdou2022wordorder,papadimitriou2022role}. Robust IR more broadly studies adversarial changes, distribution shift, and performance variance~\cite{liu2024robustnir}.

Structured inputs make serialization a particularly direct concern. TableFormer targets unwanted row- and column-order bias in table reasoning~\cite{yang2022tableformer}. Most closely related, Bhandari et al. show that semantically equivalent table serializations alter embeddings and retrieval results across retriever families, and study centroid-based representational stability~\cite{bhandari2026tabular}. We therefore do not claim to discover serialization sensitivity in neural retrieval. We extend the question to Web product discovery, where we measure item-level Top-$K$ visibility for human-judged relevant products, distinguish a direct item intervention from a catalog-wide change, quantify candidate-generation loss before reranking, and evaluate robustness jointly with relevance. Ranking-exposure research motivates attention to cutoffs~\cite{joachims2017unbiased,singh2018fairness}, while our evidence remains at the retrieval stage rather than realized user exposure.
